\scanpage{6}
\mspage{3}
\uline{Analyseur} $\Omega$: c'est un pseudo-anneau commutatif
(i.e. un groupe additif avec multiplication associative biadditive, commutative, pas néc. unitaire)
les op. notées $x+y$, $xy$
avec une loi \deleted{de \unsure{composition} \unsure{associative} (\unsure{qui est comm.}) \unsure{opp. à} loi additive}
\inserted{de composition de monoïde}
\[
(F,G)\longmapsto F\circ G
\]
unitaire (unité notée $\deleted{\unsure{I}} I$ ou $T$) pas néc. comm.

\uline{Axiomes}\quad 1) \uline{Pour mémoire}
\[
\left\{
\begin{array}{l}
+\ \mathrm{ass,commutative,\ avec}\ 0\\
xy\ \mathrm{associative,\ commut.,biadditive}\\
x\circ y\ \mathrm{associative,unitaire}
\end{array}
\right.
\]
2) $\forall G\in\Omega$, $F\longmapsto F\circ G:\Omega\longrightarrow\Omega$ est un homo. de ps.-anneaux i.e.
\[
\left\{
\begin{aligned}
(F+F')\circ G&=F\circ G+F'\circ G\\
(FF')\circ G&=(F\circ G)(F'\circ G)
\end{aligned}
\right.
\]
3) (Facultatif) a) $\forall F\in\Omega$, $\exists$ famille finie $(F_i',F_i'')$, avec $F_i',F_i''\in\Omega$, telle que
\[
F\circ(G'+G'')=\sum_i(F_i'\circ\unsure{G_i'|G'})(F_i''\circ G'')
\]
b) $\forall F\in\Omega$, $\exists$ familles finies $(\Phi_\alpha',\Phi_\alpha'')$
\[
F\circ(G'G'')=\sum_i(\Phi_\alpha'\circ G')(\Phi_\alpha''\circ G'')
\]

\uline{NB} Quid de $F\circ0$? On a $(F\circ0)\circ G=F\circ(\underbrace{0\circ G}_{0})=F\circ0$,
donc $F\circ0$ est ``une constante'', au sens suiv.

\uline{Déf} Un $\lambda\in\Omega$ est dit ``constant'' si $\lambda\circ F=\lambda$ $\forall F\in\Omega$.
\margin{[Ne dépend que de la structure $F\circ G$]}

\uline{Proposition} L'ens $\Omega_0$ des constantes de $\Omega$ est un sous-\inserted{ps.}anneau de $\Omega$
[stable par $\circ$, car $\lambda\circ\mu=\lambda$ pour $\lambda,\mu\in\Omega$).
Pour $F\in\Omega$, $\lambda\in\Omega_0$, on a $F\circ\lambda\in\Omega_0$
(car $(F\circ\lambda)\circ G=F\circ(\underbrace{\lambda\circ G}_{\lambda})=F\circ\lambda$).

Ainsi, $\Omega$ ``opère'' sur $\Omega_0$ par $(F,\lambda)\longmapsto F(\lambda)\overset{\mathrm{def}}{=}F\circ\lambda$.
\[
\Omega^0=\{F\in\Omega\mid F\circ0\overset{\mathrm{def}}{=}F(0)=0\}.
\quad\text{On a }\Omega=\Omega_0\oplus\Omega^0.
\]

\uline{Déf.} Analyseur \uline{unitaire}: la multiplication $xy$ admet une unité $1$, et celle-ci est une constante i.e.
$1\circ F=1$ $\forall F\in\Omega$.

Donc $\Omega_0$ est alors un anneau unitaire, admettant
$\unsure{\equiv}$ unité que $\Omega$
\margin{5}
